\section{Algebraizing the analytic flat frame on its actual line}
\label{sec:faithful-flat-line-extension}

The formal flat-frame construction uses three separate operations:
extend the analytic coefficient as a unit in the prescribed line,
algebraize that section, and descend its generator through the
coefficient extension and the resolution.  We give the line arguments
here in the order used by \srcid{GagaLineGeneratorComplex},
\srcid{GagaLineGeneratorEmbedding}, and
\srcid{ProperPullbackLineBundle}.  These arguments do not construct the
integral moduli model or prove its boundary growth estimates.

We use the following ordinary background form of GAGA: on a proper
smooth complex algebraic scheme, holomorphic sections of an algebraic
line bundle are precisely the analytifications of its algebraic
sections.  For the projective form see
\href{https://numdam.org/articles/10.5802/aif.59/}
{J.-P. Serre, \emph{G\'eom\'etrie alg\'ebrique et g\'eom\'etrie analytique},
Ann. Inst. Fourier 6 (1956), 1--42}.
The source calls its proper-scheme section formulation
\srcid{GagaSections}.  No assertion that an arbitrary analytic function
on an open variety is algebraic is being used.

\subsection{The line, its local generators, and the extension condition}

Let \(X_0\) be a smooth integral component of a proper smooth scheme
over a field \(K\) of characteristic zero, and let
\(\Omega=K(X_0)\).  An invertible lattice in \(\Omega\) means an
invertible subsheaf \(\mathcal L\) of the constant sheaf \(\Omega\).
Thus, locally, there is \(T\in\Omega^\times\) with
\[
                         \mathcal L_x=T\mathcal O_{X_0,x}.
 \tag{F1}\label{eq:faithful-line-generator}
\]
This is the rank-one case of the source's \(\srcid{Lat}\) and
\(\srcid{IsVB}\) conditions.  Saying that \(t\) generates at \(x\)
means \(t\ne0\) and
\(\mathcal L_x=t\mathcal O_{X_0,x}\).
The two generators \(T,t\) then satisfy
\[
             t/T\in\mathcal O_{X_0,x}^{\times}.
 \tag{F2}\label{eq:faithful-line-unit-ratio}
\]
Indeed, writing \(t=uT\) and \(T=vt\) gives \(uv=1\) after
cancelling \(T\ne0\).  Conversely, (F2) implies equality of the
two principal lattices.  Both the ratio and its inverse are needed.

Suppose \(W\subset X_0\) is a nonempty open and \(1\) generates
\(\mathcal L\) at every point of \(W\).
For an embedding \(\sigma:K\hookrightarrow\mathbb C\), let \(F\)
be a complex-valued function on the \(\sigma\)-complex points of \(W\).
The extension condition at \(x\in X_0(\mathbb C)\) is the following:
there is a rational local generator \(t\) of \(\mathcal L\) at \(x\)
such that, in every holomorphic coordinate chart through \(x\),
\(F/t\), initially defined on the inverse image of \(W\), extends
holomorphically near \(x\), with nonzero value at \(x\).
The extended function can be restricted to a smaller neighborhood
where it is everywhere nonzero.
This is \srcid{ExtendsUnitAt}; it includes a unit assertion.
A bound on the norm of \(F\), without this coefficient statement in
the specified line, would be insufficient.

Here and below a rational expression is evaluated only where it is
regular.  In particular \(t\) is a regular unit on
\(W\) near \(x\) whenever both \(t\) and \(1\) generate the line.
The source's \srcid{ValAt} notation keeps these evaluations explicit.

\begin{srlemma}[Regularity and local generators]
\label{lem:faithful-line-local-units}
Let \(T\) generate \(\mathcal L\) throughout an open \(U\).
If \(t\) generates at \(x\in U\), then, after shrinking to
\(N\subset U\) containing \(x\), there is a unit
\(v\in\Gamma(N,\mathcal O)^\times\) whose rational value is \(t/T\).
If a rational function is regular at every point of an open, it has
a unique regular section on that open; if its inverse is also regular
everywhere, this section is a unit.
\end{srlemma}
\begin{proof}
At \(x\), both \(t/T\) and \(T/t\) belong to the local ring by (F2).
Represent them by regular sections on neighborhoods of \(x\).
On a common smaller neighborhood their product is \(1\), because
it is \(1\) in the function field and the structure sheaf embeds
in that field.  These sections give \(v,v^{-1}\).
For the second assertion choose a representative of the given
rational function around every point.
On intersections the representatives have the same rational value,
so are equal by injectivity into the function field.
The sheaf axiom glues them, and proves uniqueness.
Apply the same construction to the inverse and multiply to get \(1\).
This also proves the source's transition-unit construction without
an additional gluing assumption.
\end{proof}

\begin{srlemma}[The transition units of the actual lattice]
\label{lem:faithful-line-transition-units}
Choose open sets \(U_i\) covering \(X_0\), and generators
\(T_i\in\Omega^\times\) of \(\mathcal L|_{U_i}\).
There are unique transition units \(g_{ij}\) on \(U_i\cap U_j\)
with rational values \(T_j/T_i\).  They satisfy
\[
                   g_{ii}=1,\qquad g_{ij}g_{jk}=g_{ik}.
 \tag{F3}\label{eq:faithful-line-cocycle}
\]
A collection of regular, or holomorphic, coefficients \(h_i\)
defines a section of this same line exactly when
\(h_i=g_{ij}h_j\) on every overlap.
\end{srlemma}
\begin{proof}
At every point of the overlap the ratio \(T_j/T_i\) and its inverse
are regular by (F2).  Apply the preceding lemma.
The identities (F3) hold in \(\Omega\), hence for the unique sections.
Finally \(h_iT_i=h_jT_j\) is precisely
\(h_i=(T_j/T_i)h_j\), which gives the asserted convention.
The source extends this cover by the complement of the chosen
component and sets all its coefficients to zero there.
This causes no mixed-overlap conditions: a component of a smooth
locally noetherian scheme is open and closed.
\end{proof}

\subsection{From local analytic units to an algebraic generator}

\begin{srproposition}[The complex generator, \srcid{gaga_rat_C}]
\label{prop:faithful-complex-flat-generator}
Let \(K=\mathbb C\) in the preceding setup, and assume the stated
proper GAGA section theorem.  Suppose \(F\) satisfies the extension
condition at every complex point of \(X_0\).
There is \(f\in\Omega^\times\), generating \(\mathcal L\) at every
point of \(X_0\), with
\[
                           f(y)=F(y)\quad(y\in W(\mathbb C)).
 \tag{F4}\label{eq:faithful-complex-flat-generator}
\]
\end{srproposition}
\begin{proof}
Fix \(i\) and a point \(x\in U_i(\mathbb C)\).
Take the generator \(t\) in the extension condition and the local
unit \(v=t/T_i\) of Lemma~\ref{lem:faithful-line-local-units}.
In a chart through \(x\), let \(h'\) be the holomorphic unit extending
\(F/t\).  Shrink the chart so that \(h'\) remains nonzero.
Then
\[
                              H=h'\,v
 \tag{F5}\label{eq:faithful-hchart}
\]
is a nonvanishing holomorphic function there.
On the part of the chart lying in \(W\), it satisfies
\(H\,T_i=F\).
This is exactly the construction
\srcid{exists_hChart_of_extendsUnitAt}.

Any nonempty open subset of a holomorphic coordinate chart on
\(X_0(\mathbb C)\) meets \(W(\mathbb C)\).
To justify this density, the complement of \(W\) is a proper
algebraic closed subset of the integral component.
If it contained a chart-open set, its defining regular equations
would vanish on an analytic open.  The identity theorem for
holomorphic functions would make them vanish on the component,
contrary to properness of that closed subset.
Equivalently one can use the source's
\srcid{smooth_algebra_eq_zero_of_vanish_open} chart criterion.

On two overlapping charts the constructions of \(H\) agree after
the coordinate change on their dense intersection with \(W\).
Their holomorphic continuations agree everywhere by continuity.
This proves uniqueness of the coefficient in the fixed frame
\(T_i\), so all charts glue to a holomorphic function \(h_i\)
on \(U_i(\mathbb C)\).
In the source, \srcid{HChart.chart.factor} provides the coordinate
change and \srcid{overlapping_analytic_charts_ratio} performs this
same dense-set comparison.

For different algebraic frames the dense-set identity is
\[
                      h_i=(T_j/T_i)h_j.
 \tag{F6}\label{eq:faithful-hchart-transition}
\]
The transition ratio is the regular unit of (F3).
The same chart argument extends (F6) across the complement of \(W\).
The functions \(h_i\) are holomorphic in every coordinate chart,
not merely continuous: locally they are the functions (F5)
composed with holomorphic coordinate changes.
They therefore give a holomorphic section \(s\) of
\(\mathcal L^{\mathrm{an}}\), nonvanishing at every complex point.

Proper GAGA algebraizes \(s\).
Its algebraic coefficient on each \(U_i\) has nonzero value at
every complex point, so is a unit.
For otherwise its zero locus on a finite-type complex affine
subopen would be nonempty, and a maximal ideal of its coordinate
ring would give a complex closed point where it vanishes.
Consequently the algebraic section generates the line everywhere.
Its value at the generic point is \(f\in\Omega^\times\).
Since its coefficient in \(T_i\) is \(h_i\), the identity
\(h_iT_i=F\) on \(W\) proves (F4).
The source's complex theorem permits a nonzero scalar \(\kappa\)
and supplies the particular value \(\kappa=1\).
\end{proof}

\subsection{The coefficient-by-coefficient descent used by the source}

The next lemma records the actual algebraic descent, rather than
assuming that an arbitrary complex generator is defined over \(K\).
For subrings \(S_i\subset\Omega\) containing \(K\), the notation
\(S_i\otimes_K C\) below means its image in
\(\Omega\otimes_K C\).  Membership means a finite sum of tensors
\(a\otimes c\) with \(a\in S_i\).

\begin{srlemma}[Tensor coefficient descent, \srcid{descent_core}]
\label{lem:faithful-flat-tensor-descent}
Let \(K\subset\Omega\) and \(K\subset C\) be fields and assume
\(\Omega\otimes_K C\ne0\).  Let \(S_i\subset\Omega\) contain \(K\),
and fix \(e_i\in\Omega^\times\).  Suppose
\(\tau,\upsilon\in\Omega\otimes_K C\) satisfy
\[
 \tau\upsilon=1,\qquad
 \tau=e_i x_i,\qquad e_i\upsilon=y_i,\qquad
 x_i,y_i\in S_i\otimes_K C \quad\hbox{for every }i.
 \tag{F7}\label{eq:faithful-descent-local-tensors}
\]
Assume the following two precise scalar properties:
\[
 \begin{split}
 a\ne0,\ a\in\bigcap_iS_i
     &\Longrightarrow a\in K,\\
 t\ne0,\ t/e_i,e_i/t\in S_i\ (\hbox{all }i),\
                 u/e_i\in S_i\ (\hbox{all }i)
     &\Longrightarrow u\in Kt .
 \end{split}                                      \tag{F8}
 \label{eq:faithful-descent-scalar-properties}
\]
Then there exist \(t\in\Omega^\times\), \(c\in C\), such that
\[
            t/e_i,e_i/t\in S_i\quad(\hbox{all }i),
                   \qquad \tau=t\otimes c.
 \tag{F9}\label{eq:faithful-descent-result}
\]
\end{srlemma}
\begin{proof}
Choose a \(K\)-basis \((b_j)\) of \(C\).
It gives an \(\Omega\)-basis \((1\otimes b_j)\) of
\(\Omega\otimes_K C\).  Every tensor has a finite expansion in
this basis.  If \(z\in S_i\otimes_K C\), all its coefficients
belong to \(S_i\): this holds for a simple tensor by expanding its
second factor in the \(K\)-basis, and is preserved by finite sums.

Write the coefficients of \(\tau\) as \(t_j\), and those of
\(\upsilon\) as \(u_j\).  Equation (F7) gives, for all \(i,j\),
\[
                         t_j/e_i\in S_i,\qquad
                         e_i u_j\in S_i .
 \tag{F10}\label{eq:faithful-descent-coefficients}
\]
Because \(\tau\upsilon=1\ne0\), neither tensor is zero.
Choose \(j_0,k_0\) with \(t=t_{j_0}\ne0\) and \(u=u_{k_0}\ne0\).
Multiplying the two memberships in (F10) shows that \(tu\)
belongs to every \(S_i\).  The first property in (F8) therefore
gives \(tu=\kappa\in K^\times\).
It follows that
\[
                   e_i/t=(e_i u)\kappa^{-1}\in S_i.
\]
Together with (F10), this proves that \(t\) is a unit multiple
of every \(e_i\) in the corresponding ring.

The second property in (F8) now applies to every \(t_j\),
and gives \(t_j=c_jt\) with \(c_j\in K\).
Only the finitely many \(t_j\) in the support of \(\tau\)
are needed.  With \(c=\sum_jc_jb_j\in C\), bilinearity gives
\(\tau=t\otimes c\).
This proves every assertion of (F9), including the inverse
membership.  The proof does not assume that the chosen basis
of \(C/K\) is finite.
\end{proof}

For a proper integral \(K\)-scheme with \(K\) algebraically closed
in its function field \(\Omega\), a rational function regular
everywhere is in \(K\).
One proof, matching the source's regular-everywhere step, is as follows.
A regular rational function is a global section by the sheaf
argument above.  The algebra of global functions on a proper
finite-type scheme over \(K\) is finite dimensional over \(K\).
Multiplication by that section is a linear operator on this
nonzero finite-dimensional algebra.
Its characteristic polynomial annihilates the operator and,
on applying to \(1\), annihilates the section.
Thus the section is algebraic over \(K\), and the assumed
constant-field condition puts it in \(K\).
This proves the first property of (F8) for
\(S_i=\mathcal O_{X_0,x_i}\), indexed by all points.
The second follows by applying this to \(u/t\), because
\((u/e_i)(e_i/t)\) is regular everywhere.
The case \(u=0\) gives the scalar \(0\).
This is the field-base specialization of
\srcid{LineLatticeGenericSections.mem_lat_genFibre_iff_const_mul}.

\begin{srproposition}[Descent of the generator]
\label{prop:faithful-flat-generator-descent}
Let \(X_0\) be proper integral over a characteristic-zero field \(K\),
with function field \(\Omega\), and assume \(K\) is algebraically
closed in \(\Omega\).
Let \(C/K\) be a field extension such that \(\Omega\otimes_K C\)
is a domain, and let \(\Omega'\) be its fraction field.
Suppose \(t'\in\Omega'^\times\) generates the pullback of an
invertible lattice \(\mathcal L\subset\Omega\) everywhere on
\(X_0\times_K C\).
Then there are a generator \(t\in\Omega^\times\) of
\(\mathcal L\) everywhere on \(X_0\), and \(c\in C^\times\),
such that \(t'=ct\) in \(\Omega'\).
\end{srproposition}
\begin{proof}
Choose affine neighborhoods \(U_i\) with rational generators \(e_i\)
of \(\mathcal L\); their rings embed in \(\Omega\).
Their base changes are affine, with rings
\(\Gamma(U_i,\mathcal O)\otimes_K C\), embedded in
\(\Omega\otimes_K C\) by flatness over the field \(K\).
Since both \(t'\) and \(e_i\) generate the pulled line,
\(t'/e_i\) and \(e_i/t'\) are regular everywhere on \(U_{i,C}\).
The sheaf argument makes them elements of its coordinate ring.
Taking any one nonempty \(U_i\) shows that \(t'\) and \(t'^{-1}\)
actually lie in \(\Omega\otimes_K C\), despite their initially
being specified in its fraction field.
Call these tensors \(\tau,\upsilon\).
Their product is \(1\), since their images in the fraction field
have product \(1\).

At each point \(x\) take a local generator \(e_x\) defined on an
affine neighborhood and set \(S_x=\mathcal O_{X_0,x}\).
The preceding affine calculation gives (F7) in
\(S_x\otimes_K C\).  The proper constant-field argument gives
(F8).  Lemma~\ref{lem:faithful-flat-tensor-descent} produces
\(t\) with \(t/e_x,e_x/t\in S_x\), hence a generator everywhere,
and \(t'=ct\).
Since \(t'\) and \(t\) are nonzero, \(c\ne0\).
This is the coefficient argument in
\srcid{ProperModelGenericFibreCompare.exists_gen_of_extension};
the more general source statement allows a domain base
and restricts to its generic fibre.
\end{proof}

For the application \(C=\mathbb C\) with its fixed embedding
\(\sigma:K\hookrightarrow\mathbb C\), the source constructs
\(\Omega'=\operatorname{Frac}(\Omega\otimes_{K,\sigma}\mathbb C)\).
In characteristic zero the constant-field condition makes
\(\Omega/K\) regular; the source uses
\srcid{alg_closed_tensor_isDomain} for the asserted domain
property.  The maps of generic points and the stalk maps commute
with this base change.
Local generators, their unit ratios, and the analytic extension
condition pull back: the ratio is evaluated through the same
stalk homomorphism, and uniqueness of a rational function's
regular value identifies the two evaluations.
Thus Proposition~\ref{prop:faithful-complex-flat-generator}
applies upstairs.
If it gives \(f'\) with \(f'(y)=\kappa F(y)\), and
\(f'=ct\) is the descent just proved, then
\[
                           t(y)=(\kappa/c)F(y).
 \tag{F11}\label{eq:faithful-flat-generator-scalar}
\]
Both scalars are nonzero.  This recovers exactly the conclusion
of \srcid{GagaLineGeneratorEmbedding.gaga_rat}.

\subsection{Descent through the proper resolution}

\begin{srlemma}[The normal proper pushdown of a generator]
\label{lem:faithful-flat-proper-pushdown}
Let \(X_0\) be normal integral and locally noetherian,
let \(f:X'\to X_0\) be proper and birational on the chosen
generic component, and identify their function fields \(\Omega\).
Let \(\mathcal L\subset\Omega\) be an invertible lattice in a
neighborhood of \(x\in X_0\).
Suppose \(t\in\Omega\) generates \(f^*\mathcal L\) at every
point of the chosen component of \(X'\) whose image is a
generization of \(x\).
Then \(t\) generates \(\mathcal L\) at \(x\).
\end{srlemma}
\begin{proof}
The generic point is among the points in the hypothesis, so \(t\ne0\).
Shrink to a neighborhood of \(x\) on which \(e\ne0\) generates
\(\mathcal L\), and set \(r=t/e\).
It suffices to prove \(r,r^{-1}\in\mathcal O_{X_0,x}\).

Let \(q\) be any codimension-one point generizing \(x\).
Normality and local noetherianity make
\(V=\mathcal O_{X_0,q}\) a discrete valuation ring with fraction
field \(\Omega\).
The chosen generic point of \(X'\), together with
\(\operatorname{Spec}V\to X_0\), gives the valuative diagram
for the proper map \(f\).
The valuative criterion supplies a lift
\(\operatorname{Spec}V\to X'\); its closed point \(q'\) lies
over \(q\) on the chosen component.
By hypothesis \(t\) generates the pullback line there,
while \(e\) also generates there by pullback.
Consequently \(r,r^{-1}\) belong to
\(\mathcal O_{X',q'}\), and the lift maps them into \(V\).
The induced map on fraction fields is the fixed identity,
so these are the original rational functions \(r,r^{-1}\).
Thus both belong to every such height-one localization.

A noetherian normal local domain is the intersection of its
height-one localizations in its fraction field.
Apply this intersection statement to
\(\mathcal O_{X_0,x}\).
It gives \(r,r^{-1}\in\mathcal O_{X_0,x}\), as required.
If this local ring has dimension zero it is the function
field and the conclusion is immediate.
This is the source proof: its \srcid{exists_over} supplies
\(q'\), and \srcid{line_lattice_generator_codim_one}
makes the last intersection step.
No arbitrary specialization of a quotient is used.
\end{proof}

\subsection{Where these arguments enter the geometric construction}

In \srcid{EigenFramesAnalytic} one first takes the nonempty generic
open where the chosen eigenlattice bases are frames.
The line to which the preceding results apply is the actual
wedge-ratio lattice, formed from the two conjugate eigenspaces;
its generators are ratios of top exterior products of those frames.
The source's \srcid{ExistsResolutionGenericFibre} supplies a proper
smooth resolution, an isomorphism on this open, and a simple
normal-crossing complement.  The pullback wedge-ratio lattice is
locally free of rank one there, including along that complement.

The construction-specific local assertion is
\srcid{EigenFrameFunctionCompare}:
the analytic flat function \(\Phi\), after expression in any
local wedge frame, extends as a nonvanishing holomorphic function.
Its producer \srcid{AffineChartBundleSections} uses the actual
boundary chart, the actual eigenvectors and their pullbacks,
and \srcid{BoundaryChartFrameGrowth}.
Once this precise local assertion is proved, the steps printed
above algebraize its section and descend the generator through
the proper resolution to the normal generic-fibre model.
\srcid{WedgeEigenLineGenerator} then turns that generator into the
frame function \(\varphi\).
The comparison \srcid{FrameFunctionsFlatAgree} is made on this
same line: two such generators have ratio regular and invertible
everywhere, hence a scalar in the coefficient field by the
constant-field argument.  Determinants of basis changes transfer
the equality to every adapted frame.

For clarity, none of the preceding line arguments supplies
\srcid{Model.N3}, its polarized uniformization or torsion pairing,
or the boundary-chart unit assertion just identified.
Nor do they identify a generic eigenspace with a free integral
eigenmodule at a ramified special point.
Those are separate construction-specific premises.
The completed results of this section are the transition and
analytic-unit gluing, algebraization, coefficient descent, and
normal proper pushdown, with their precise hypotheses.
Their use does not conceal one of those remaining model inputs.
